\documentclass[a4paper,12pt]{article}
\usepackage[utf8]{inputenc}
\usepackage[T1]{fontenc}
\usepackage[margin=2.5cm]{geometry}
\usepackage{ctex} % 支持中文
\usepackage{setspace} % 设置行距
\usepackage{titlesec} % 标题格式
\usepackage{xcolor} % 颜色支持
\usepackage{enumitem} % 列表优化

% 设置行距为1.5倍，增加正式感
\onehalfspacing

% 设置标题格式
\titleformat{\section}{\Large\bfseries\color{blue!40!black}}{}{0.5em}{}
\titleformat{\subsection}{\large\bfseries\color{gray!80!black}}{}{0.5em}{}

% 列表设置
\setlist[itemize]{leftmargin=*}

% 元数据
\title{\textbf{2024年度环境、社会及管治（ESG）报告节选}}
\author{公司董事会办公室}
\date{2024年3月}

\begin{document}

\maketitle

% 依据指令，已移除摘要和目录

\section{1. 前言：战略协同与可持续发展框架综述}

2024年，被界定为海南自由贸易港封关运作准备工作的攻坚之年，亦是本公司（以下简称“公司”）作为区域性基础设施建设与临空产业运营主体，深化高质量发展路径的关键时期。鉴于公司在“一主一次，两支一货运”区域机场群协同发展战略中的核心地位，以及其对于国家对外开放战略布局的实质性支撑作用，本报告旨在系统性地披露公司在环境、社会及管治（ESG）维度的管理效能与绩效表现。公司致力于通过建立健全的内部控制体系与可持续发展治理架构，确保“外防输入、内保畅通”运营目标的达成，并在此基础上，通过对环境足迹的精细化管控及社会价值的多元化创造，为包括股东在内的所有利益相关方构建长期稳健的价值回报机制。

\section{2. 环境范畴（Environmental）：碳排放管控与资源集约化利用}

在应对全球气候变化及响应国家“双碳”战略目标的宏观背景下，公司已将环境管理纳入企业核心战略范畴。报告期内，公司依据适用的法律法规及行业标准，实施了严格的温室气体排放核算与能源审计程序，旨在通过技术迭代与管理优化，实现运营碳强度的持续下降。

\subsection{2.1 温室气体排放与资源消耗量化绩效}

依据ISO 14064标准及相关行业温室气体核算指南，公司对2024年度的碳排放数据进行了系统性盘查。经核算，报告期内公司温室气体排放总量录得 \textbf{51,518.71} 吨二氧化碳当量（tCO$_2$e）。对此数据进行结构性分析可知，源自固定源与移动源燃烧的直接排放（范围1）为 \textbf{11,214.39} tCO$_2$e，而源自外购电力消费的间接排放（范围2）则高达 \textbf{40,304.32} tCO$_2$e。基于年度营收规模测算，单位营收碳排放强度为 \textbf{0.12} 吨二氧化碳当量/万元。

在能源与资源消耗的审计方面，2024年度综合能源消耗量折合标准煤 \textbf{13,538.22} 吨。具体而言，外购电力消耗量为 7,511.05 万千瓦时，移动源汽油消耗量为 19.10 万升。关于水资源管理与废弃物处置，报告期内水资源消耗总量为 145.94 万吨，其中通过再生水利用系统实现的中水回用量为 13.61 万吨；同期产生的无害废弃物与有害废弃物分别为 6,230.58 吨与 0.58 吨，且上述废弃物的合规处置率均达到了 100\%。

\subsection{2.2 绿色机场建设路径与减排技术应用}

为达成“绿色机场”的建设目标，公司采取了一系列旨在降低航空器地面运行能耗及提升可再生能源利用率的技术干预措施。

关于辅助动力装置（APU）替代设施的部署，该技术旨在通过地面电源装置（GPU）及地面空调组件（PCA）的介入，替代航空器在地面停泊期间对机载APU的依赖，从而显著削减航空燃油消耗及相应的温室气体与噪声排放。截至报告期末，旗下核心门户机场已投入运营的APU替代设施共计 \textbf{31套}（涵盖18个近机位与13个远机位）。经测算，该项目的实施在2024年度实现了约 \textbf{25,382} 吨二氧化碳当量的减排效应。基于上述成效，旗下核心门户机场被授予“双碳机场”二星级认证。

在可再生能源系统的整合应用方面，旗下核心商业综合体实施了屋顶分布式光伏项目，其年均发电量达 450 万千瓦时，占该建筑体总电力需求的约 10\%，折算年度碳减排量为 2,263.59 吨；与此同时，旗下临空产业园利用闲置屋顶资源，完成了装机容量为 8.6MWp 的光伏发电系统建设，预计年均碳减排贡献可达 25.75 万吨。此外，在场内车辆电动化转型方面，截至2024年末，核心机场区域内的新能源车辆保有量已提升至 221 辆（占比 43.6\%），同期配套建设充电桩设施约 100 个，该项措施在报告期内实现了 761 吨的碳排放削减。

\subsection{2.3 绿色建筑认证体系}

在建筑资产的开发与运营全生命周期中，公司严格遵循国际及国内绿色建筑评价标准。省会城市在建地标性项目“海南中心”（双子塔）已成功获取 \textbf{LEED CS 金级预认证}，体现了其在建筑节能与环境设计领域的卓越表现；位于某新区的F07项目亦经评估符合并获得了绿色建筑一星级标准认证。

\section{3. 社会范畴（Social）：运营韧性构建与利益相关方价值共创}

公司在确保航空安全运营底线的基础上，致力于通过服务流程的再造与应急响应机制的强化，提升运营韧性；同时，坚持多元包容的人力资源政策，并积极履行企业公民责任。

\subsection{3.1 核心运营指标与业务恢复情况}

报告期内，受宏观经济复苏及旅游市场需求反弹的驱动，公司主要运营指标呈现稳健增长态势。旗下控股及管理输出的9家机场年度总体起降架次达 \textbf{16.11 万架次}，旅客吞吐量共计 \textbf{2,462.62 万人次}。值得注意的是，旗下核心门户机场旅客吞吐量重回“2000万级”梯队（位列全国第29位），且重点区域机场成功开通了至吉隆坡的首条国际客运航线，标志着区域航空网络的进一步拓展。

\subsection{3.2 安全管理体系（SMS）与极端天气应急响应}

公司建立并维持了一套严谨的安全管理体系，并针对极端气象灾害制定了专项应急预案，以确保业务连续性。

针对2024年9月超强台风“摩羯”的应急响应与处置，公司启动了“防汛防风专项应急预案”，运管委实施了包括航班动态调整及航空器系留加固在内的预防性措施。台风过境后，迅速组织设施修复与适航检查，9月3日至7日期间，核心机场在确保安全的前提下迅速恢复运行，累计保障航班 \textbf{1,300} 架次，有效保障了区域交通枢纽的畅通。

在数字化安全管控平台的应用方面，公司上线了“安全智控”平台，以实现安全管理要素的数字化与流程化。该平台关联安全制度手册条款 9,626 项，全年流转并处置电子检查单 31 万份。基于上述管理措施，公司本年度实现了一般征候万架次率为 0、责任事故发生率为 0 的安全绩效，且安全培训全员覆盖率达到 100\%。

\subsection{3.3 服务流程优化与创新}

为提升旅客通关效率与入境体验，公司实施了查验流程集约化与入境服务便利化措施。具体而言，实施“一机双屏”海关与安检融合查验模式，通过图像共享技术实现旅客“一次过检”；并在核心机场设立“一站式综合服务中心”，集成支付绑定、通信服务及文旅咨询功能，以解决外籍旅客入境服务的“最后一公里”问题。

\subsection{3.4 人力资本与社会投入}

截至2024年末，公司员工总数为 10,611 人。在人力资源构成的多元化方面，女性员工占比为 38.9\%（4,133人），少数民族员工人数为 179 人。公司在员工培训与发展方面的年度投入总额为 1,119.54 万元，人均培训时长达 105.88 小时。在社会公益领域，公司投入乡村振兴资金 190 万元，年度对外捐赠资金 15.48 万元。此外，员工志愿服务参与人次超过 10,000 人次，累计总时长达 500,000 小时，并圆满完成了博鳌亚洲论坛年会等重大活动的保障任务。

\section{4. 管治范畴（Governance）：治理架构与合规监督体系}

公司致力于构建符合上市公司监管要求的高标准治理架构，通过强化内部控制与合规管理，保障决策的科学性与透明度。

\subsection{4.1 治理结构与合规监管}

公司构建了规范的“五会一层”治理体系。董事会由9名成员组成，其中包含3名独立董事，符合独立性监管要求，且女性董事占比为 22\%。报告期内，共召开股东大会5次，董事会会议13次，并发布各类公告70份，严格履行了信息披露义务。

在风险管控与商业道德方面，公司实施了“三全三化”大监督体系，旨在实现对关键业务流程与风险点的全覆盖监控。全年共识别风险点 358 个并制定了相应的管控措施。在反腐败培训方面，管理层与员工层的覆盖率分别达到 95\% 与 90\%，累计培训时长 400 小时。此外，供应商廉洁承诺书签署率达到 100\%，并对23家工程类供应商进行了严格遴选。

\subsection{4.2 间接经济影响}

作为自贸港建设的积极参与者，公司通过参股投资方式布局离岛免税业务，对区域经济产生了实质性的间接经济影响。2024年度，公司间接参与的5家离岛免税店实现线下销售额约为 \textbf{52 亿元}，占全岛市场份额约 \textbf{17\%}，有效促进了消费回流与地方税收增长。

\end{document}